\bigskip
\bigskip
\subsubsection*{Խնդիրներ և հարցեր թեմա 16-ի վերաբերյալ}

\begin{enumerate}[label=\thesection.\arabic*.]

% 16.1
\item Ճի՞շտ է պնդումը. Ցանկացած անտիդիսկրետ տարածություն գծային կապակցված է:

% 16.2
\item Ապացուցեք. Գծային կապակցված ամեն մի $X$ տարածության ցանկացած $X/R$ ֆակտոր-տարածություն գծային կապակցված է:

% 16.3
\item Ճի՞շտ է պնդումը. Եթե տոպոլոգիական տարածությունների $X \times Y$ արտադրյալը գծային կապակցված է, ապա $X$-ը և $Y$-ը գծային կապակցված են:
\begin{hint}
Դիտարկեք $X \times Y$ արտադրյալի $P_X$ և $P_Y$ կանոնական պրոյեկցիաները $X$-ի և $Y$-ի վրա և օգտվեք թեորեմ 5-ի ապացույցից:
\end{hint}

% 16.4
\item Ապացուցեք հետևյալ թեորեմը. $X \times Y$ արտադրյալը գծային կապակցված է այն և միայն այն դեպքում, երբ գծային կապակցված են $X$-ը և $Y$-ը:
\begin{hint}
Պայմանի բավարարությունը ապացուցելու համար դիտարկենք $X \times Y$-ի որևէ երկու $(x_1, y_1)$ և $(x_2, y_2)$ կետեր: Դիցուք $f_1: I \to X$ ուղին միացնում է $x_1$ կետը $x_2$ կետին, իսկ $f_2: I \to Y$ ուղին միացնում է $y_1$ կետը $y_2$ կետին: Օգտվելով թեմա 14-ում թեորեմ 3-ից դիտարկեք $(f_1, f_2): I \to X \times Y; (f_1, f_2)(t) = (f_1(t), f_2(t))$ արտապատկերումը:
\end{hint}

% 16.5
\item Դիցուք $X$-ը և $Y$-ը հոմեոմորֆ տարածություններ են: Ապացուցեք. $X$-ը գծային կապակցված է այն և միայն այն դեպքում, երբ գծային կապակցված է $Y$-ը:

% 16.6
\item Ապացուցեք. $\mathbb{R}^n$ էվկլիդեսյան տարածության $B^n = \{x; \|x\| \le 1\}$ միավոր գունդը գծային կապակցված ենթատարածություն է:
\begin{hint}
Հիմնավորեք, որ $B^n$-ը ուռուցիկ ենթաբազմություն է:
\end{hint}

% 16.7
\item Ապացուցեք, որ $\mathbb{R}^{n+1}$ էվկլիդեսյան տարածության $S^n_+$ և $S^n_-$ միավոր կիսասֆերաները, որտեղ$S^n_+ = \{x \in \mathbb{R}^{n+1}; \|x\|=1, x_{n+1} \ge 0\}$, $S^n_- = \{x \in \mathbb{R}^{n+1}; \|x\|=1, x_{n+1} \le 0\}$գծային կապակցված ենթատարածություններ են:
\begin{hint}
Հիմնավորելով, որ $\mathbb{R}^{n+1}$-ի պրոյեկցիան $\mathbb{R}^n$-ի վրա սահմանված $p(x_1, x_2, \dots, x_{n+1}) = (x_1, x_2, \dots, x_n)$ բանաձևով արտապատկերում է $S^n_+$ և $S^n_-$ կիսասֆերաներից յուրաքանչյուրը հոմեոմորֆիզմով $B^n$ գնդին, օգտվել նախորդ երկու խնդիրներից:
\end{hint}

% 16.8
\item Ապացուցեք. $\mathbb{R}^{n+1}$ էվկլիդեսյան տարածության $S^n = \{x; \|x\|=1\}$ միավոր սֆերան գծային կապակցված տարածություն է:
\begin{hint}
Ներկայացնելով $S^n$-ը որպես $S^n_+$ և $S^n_-$ կիսասֆերաների միավորում, օգտվել նախորդ խնդրից և թեորեմ 3-ից:
\end{hint}

% 16.9
\item Ապացուցեք. $n \ge 0$ դեպքում $\mathbb{R}P^n$ իրական պրոյեկտիվ տարածությունը գծային կապակցված է:
\begin{hint}
Երբ $n=0$, $\mathbb{R}P^0$-ն բաղկացած է մի կետից: Երբ $n > 0$, ներկայացրեք $\mathbb{R}P^n$-ը որպես $S^n$ սֆերայի, կամ $S^n_+$ կիսասֆերայի ֆակտոր-տարածություն և օգտվել նախորդ խնդիրներից:
\end{hint}

% 16.10
\item Դիտարկենք $\mathbb{R}^2$ հարթության $A = \{(x, y); x > 0, y = \sin \frac{1}{x}\}$ և $B = \{(0, y); -1 \le y \le 1\}$ ենթաբազմությունները: Ապացուցեք. $A \cup B$ միավորումը կապակցված է, բայց գծային կապակցված չէ:
\begin{hint}
Թեորեմ 7-ի նմանությամբ ցույց տվեք, որ $B$-ի ոչ մի կետ հնարավոր չէ ուղիով միացնել $A$-ի որևէ կետի հետ:
\end{hint}

\end{enumerate}
